% Generated draft: requires mathematical and source review.
\documentclass[11pt]{article}
\usepackage[margin=25mm]{geometry}
\usepackage{fontspec}
\setmainfont{Latin Modern Roman}
\usepackage{amsmath,amssymb,mathtools,unicode-math}
\setmathfont{Latin Modern Math}
\usepackage{enumitem,needspace,xurl,hyperref}
\setlength{\parindent}{0pt}
\setlength{\parskip}{5pt}
\setlength{\emergencystretch}{3em}
\newcommand{\N}{\mathbb{N}}
\newcommand{\Z}{\mathbb{Z}}
\newcommand{\Q}{\mathbb{Q}}
\newcommand{\R}{\mathbb{R}}
\newcommand{\C}{\mathbb{C}}
\newcommand{\F}{\mathbb{F}}
\newcommand{\A}{\mathbb{A}}
\newcommand{\PP}{\mathbb{P}}
\newcommand{\E}{\mathbb{E}}
\newcommand{\eps}{\varepsilon}
\newcommand{\dd}{\,\mathrm{d}}
\begin{document}
\section*{UnsolvedMath ID: 11000129}
\textbf{Source problem code:} catalogue entry
\subsection{Definitions and mathematical statement}
Retain \(B\mathcal C_2\) as the classifying space appearing in the source, interpreted by the supplied summary as a surface-category classifying space. The category \(\mathcal C_2\), its topology, and its objects and morphisms are not defined in the excerpt. Retain \(TC(\mathrm{pt})\) as its topological cyclic homology of a point; whether this denotes a spectrum, a chosen associated space, or a completed version is not specified.

\emph{Question 4.3.} Is there a geometric map from \(B\mathcal C_2\) into \(TC(\mathrm{pt})\)? A precise construction must identify the target space or the suspension-spectrum convention, state how surface-category data produce the map, and specify any required nontriviality, normalization, or compatibility. The adjective ``geometric'' is essential: an arbitrary constant map does not express the intended request, yet no formal criterion distinguishing the desired map is given. No existence of an arbitrary continuous map is substituted for the problem.

The source suggests Kevin Costello's manuscript [3], relating conformal field theory to linear Hochschild homology, as a possible starting point. It says that manuscript announces a description of the Deligne--Mumford compactification from the open moduli space, and mentions preliminary work of S\o ren Galatius and Yasha Eliashberg giving a homotopical description of a partial compactification of stable moduli space. These are attributed historical directions rather than a supplied construction of the requested map.

The summary associates \(B\mathcal C_2\) with \(\Z\times B\Gamma_\infty^+\) through Madsen--Tillmann--Weiss/GMTW and mentions an \(\Omega^\infty_{-1}\mathbb{CP}^{\infty}_{-1}\) factor of \(TC(\mathrm{pt})^{\wedge}_p\). The category, loop shifts, Thom-spectrum notation, prime, and completion conventions are not supplied. In particular, an equivalence for a loop space must not be identified with one for the unlooped classifying space. The compressed factor notation and the stated identification are retained as unresolved summary data, not asserted as a verified equivalence determining the map.
\subsection{Short English statement}
Construct the intended geometric map from the surface-category classifying space to topological cyclic homology of a point. The category, target spectrum/space convention, loop shifts, and meaning of geometric are missing; the source offers conformal-field-theory and compactification programs as possible approaches.
\subsection{Sources}
\begin{itemize}
\item Local UnsolvedMath registry, permanent ID 11000129, source code catalogue entry, Question 4.3, PDF page 105; stored title ``Question 4.3 --- .'' No individual proposer is supplied.
\item Farb, \emph{Problems on Mapping Class Groups and Related Topics}. Supplied URL: \url{https://www.math.uchicago.edu/~farb/papers/mcgbook.pdf}; recorded access 2026-07-29; extraction pdf-font-structure. The volume contains 329 explicitly typeset headings across chapters, with repeated numbering disambiguated by page.
\item The statement cites Kevin Costello [3] and preliminary work by S\o ren Galatius and Yasha Eliashberg, without full entries. The summary identifies Costello, ``Topological conformal field theories and Calabi--Yau categories,'' \emph{Adv. Math.} (2007), and refers to Madsen--Tillmann--Weiss/GMTW without detailed citations.
\item Supplied status open and OPEN-TRIAGE. The summary reports no geometric map or surface-category interpretation of the indicated completed-TC factor found in its audit. This is an editorial search assessment, not an independently verified absence claim. Current primary-literature status and redistribution terms for the publicly accessible author/publisher copy remain NEEDS\_REVIEW.
\end{itemize}
\end{document}
